%% FullCourse_10Lecs -- Lecture 9, Part 9.3b
%% Assembled by tools/lec09_assemble.py from reorg_manifest.yaml: Phase A of
%% Lec09_Reorg_Plan_2026-09-12.md as amended by Lec09_Reorg_Plan_Review_2026-09-12.md.
%% Each "% ---- <ID> · <TAG> · TIER" marker names a frame's source deck and line;
%% keep the markers until Phase B is done (lec09_assemble.py check reads them).
%% Owns: dials, ceilings, isolation, wirings, the router myth, dispatch evidence
%% Owns: the three rungs and their surfaces (/agents, --agent, --agents, -p, --bg) with one verified line per other harness
%% Owns: four-digit runs (Navier-Stokes as a dated sidebar), the ledger, the corrections table
%% Defers: 13k-call numbers -> Lec10 T6b; reliability under repetition -> T4; four mistakes at scale -> T5c

%% Shared preamble for the FullCourse_10Lecs topic decks.
%%
%% Each topic .tex \input's this file, then defines its own \title, \date
%% override (if any), \graphicspath, and \begin{document}.
%%
%% Reconstructed verbatim from the inlined copy preserved in
%% Lec10_Case_Studies/Lec10_T8_Case6_DeepRamsey.tex (lines 23-190), which is
%% explicitly annotated there as "from ../shared_preamble.tex, copied verbatim".
%% Base preamble matches MiniCourse_8hr/Slides/shared_preamble.tex; this version
%% additionally provides \sectiondivider, the conceptbox/cautionbox/examplebox/
%% takeawaybox tcolorboxes, and \compacttables used across the split decks.

\documentclass[10pt,english,aspectratio=169]{beamer}

\usepackage{ae,aecompl}
\usepackage[T1]{fontenc}
\usepackage{textcomp}
\usepackage[utf8]{inputenc}
\usepackage{babel}
\usepackage{datetime2}

\setcounter{secnumdepth}{3}
\setcounter{tocdepth}{3}

\usepackage{amsmath, amssymb, amsthm, graphicx}
\usepackage{booktabs, multirow}
\usepackage{tikz}
\usepackage{pgfplots}
\pgfplotsset{compat=1.16}
\usepackage[most]{tcolorbox}
\tcbuselibrary{skins, breakable, theorems}
\usepackage{listings}
\usepackage{xcolor}

\lstset{
  basicstyle=\ttfamily\small,
  keywordstyle=\color{blue!80!black}\bfseries,
  commentstyle=\color{green!50!black}\itshape,
  stringstyle=\color{orange!80!black},
  backgroundcolor=\color{gray!8},
  frame=single,
  rulecolor=\color{gray!40},
  breaklines=true,
  showstringspaces=false,
  numbers=none,
}

\lstdefinestyle{pythonstyle}{
    language=Python,
    basicstyle=\ttfamily\footnotesize,
    keywordstyle=\color{blue!80!black}\bfseries,
    commentstyle=\color{green!50!black}\itshape,
    stringstyle=\color{orange!80!black},
    frame=single,
    breaklines=true,
    showstringspaces=false,
    numbers=left,
    numbersep=5pt,
    numberstyle=\tiny\color{gray},
    backgroundcolor=\color{gray!8},
    rulecolor=\color{gray!40}
}

\ifx\hypersetup\undefined
  \AtBeginDocument{%
    \hypersetup{bookmarks=false,breaklinks=true,pdfborder={0 0 0},colorlinks=true,
      linkcolor=DarkText,urlcolor=RedTitle}%
  }
\else
  \hypersetup{bookmarks=false,breaklinks=true,pdfborder={0 0 0},colorlinks=true,
    linkcolor=DarkText,urlcolor=RedTitle}%
\fi

\makeatletter
\newcommand\makebeamertitle{%
  \begin{frame}[plain]
    \vspace*{1.0cm}
    \centering
    {\usebeamerfont{title}\usebeamercolor[fg]{title}\inserttitle\par}
    \vspace{1.0cm}
    {\usebeamerfont{author}\usebeamercolor[fg]{author}\insertauthor\par}
    \vspace{0.2cm}
    {\usebeamerfont{date}\usebeamercolor[fg]{date}\insertdate\par}
  \end{frame}}%
\AtBeginDocument{%
  \let\origtableofcontents=\tableofcontents
  \def\tableofcontents{\@ifnextchar[{\origtableofcontents}{\gobbletableofcontents}}
  \def\gobbletableofcontents#1{\origtableofcontents}%
}

\usetheme{metropolis}
\usefonttheme{professionalfonts}

\definecolor{LightGray}{RGB}{230,230,230}
\definecolor{RedTitle}{RGB}{0,0,200}
\definecolor{VeryLightGray}{RGB}{245,245,245}
\definecolor{DarkText}{RGB}{0,0,0}

\setbeamercolor{background canvas}{bg=white}
\setbeamercolor{title}{fg=RedTitle}
\setbeamercolor{author}{fg=DarkText}
\setbeamercolor{date}{fg=DarkText}
\setbeamercolor{frametitle}{bg=VeryLightGray,fg=DarkText}
\setbeamercolor{normal text}{fg=DarkText}
\setbeamerfont{title}{size=\large}

\setbeamersize{text margin left=5mm, text margin right=5mm}
\addtobeamertemplate{frametitle}{}{\vspace{-0.5em}}
\newcommand{\reducevspace}{\vspace{-0.5cm}}

% Shared section divider and box styles for split topic decks.
\newcommand{\sectiondivider}[2]{%
\begin{frame}[plain,noframenumbering]
\begin{center}
\vfill
{\Large\textbf{#1}\par}
\vspace{0.35cm}
{\normalsize #2\par}
\vfill
\end{center}
\end{frame}
}

\newtcolorbox{conceptbox}[1][]{
  colback=blue!3,
  colframe=RedTitle!55,
  boxrule=0.35mm,
  arc=1mm,
  left=2mm,
  right=2mm,
  top=1mm,
  bottom=1mm,
  #1
}
\newtcolorbox{cautionbox}[1][]{
  colback=red!3,
  colframe=red!50!black,
  boxrule=0.35mm,
  arc=1mm,
  left=2mm,
  right=2mm,
  top=1mm,
  bottom=1mm,
  #1
}
\newtcolorbox{examplebox}[1][]{
  colback=green!3,
  colframe=green!45!black,
  boxrule=0.35mm,
  arc=1mm,
  left=2mm,
  right=2mm,
  top=1mm,
  bottom=1mm,
  #1
}
\newtcolorbox{takeawaybox}[1][]{
  colback=VeryLightGray,
  colframe=RedTitle!55,
  boxrule=0.35mm,
  arc=1mm,
  left=2mm,
  right=2mm,
  top=1mm,
  bottom=1mm,
  #1
}
\newcommand{\compacttables}{\renewcommand{\arraystretch}{1.15}}

\setbeamertemplate{navigation symbols}{}
\setbeamertemplate{footline}{%
  \hfill%
  \usebeamercolor[fg]{page number in head/foot}%
  \usebeamerfont{page number in head/foot}%
  \insertframenumber\,/\,\inserttotalframenumber\kern1em\vskip2pt%
}

\makeatother

\author{Zhigang Feng}
% "date with time": compile timestamp so each build is identifiable.
\date{\today\ \DTMcurrenttime}


% Suppress redundant auto-generated metropolis section page; \sectiondivider frames are the dividers we keep.
\metroset{sectionpage=none}

\graphicspath{{./pic/}{./figures/}{../}}

\usetikzlibrary{positioning,arrows.meta,shapes,calc,fit,backgrounds}

%% Lec09 shared MAP slide, v2 (September 2026 reorganization): five parts,
%% eleven decks.  v1 (the July "five acts" map) is archived as
%% _archive/five_acts_2026-07/lec09_map_v1.tex.
%%
%% Input this file AFTER ../shared_preamble.tex (needs RedTitle, VeryLightGray,
%% DarkText) and call \LecNineMap{part}{sub} inside a frame:
%%   part in {1..5} highlights that part ("you are here") and, when sub names a
%%   sub-deck letter, that deck's chip -- \LecNineMap{2}{b} is T2b, \LecNineMap{4}{}
%%   is T4;  part = 0 renders the closing reprise with every part complete (the
%%   "Your Job Description" frame in T5d).
%% Each part carries its student question and its economic twin (the delegation-
%% economics thread of Lec09_Reorg_Plan_Review_2026-09-12.md, A3).
%% Filename deliberately does not match the build glob Lec*_T*.tex, so the build
%% script never tries to compile it standalone.

% Highlight chip #1 when the requested deck key equals #2 (e.g. 2b).
\newcommand{\lecmapchip}[2]{%
  \edef\lecmaptmp{#2}%
  \ifx\lecmapkey\lecmaptmp\tikzset{#1/.style={lecchipcur}}\fi}

\newcommand{\LecNineMap}[2]{%
% TikZ option lists are comma-split before TeX conditionals run, so every
% \ifnum/\ifx is resolved here, into named styles, before the picture starts.
\edef\lecmapkey{#1#2}%
\tikzset{lecparta/.style={}, lecpartb/.style={}, lecpartc/.style={},
         lecpartd/.style={}, lecparte/.style={},
         chipTone/.style={}, chipTtwoa/.style={}, chipTtwob/.style={},
         chipTthreea/.style={}, chipTthreeb/.style={}, chipTfour/.style={},
         chipTfivea/.style={}, chipTfiveb/.style={}, chipTfivec/.style={},
         chipTfived/.style={}}%
\ifnum#1=0
  \tikzset{lecparta/.style={lecpartdone}, lecpartb/.style={lecpartdone},
           lecpartc/.style={lecpartdone}, lecpartd/.style={lecpartdone},
           lecparte/.style={lecpartdone}, lecchip/.style={lecchipbase, lecchipdone}}%
\else
  \tikzset{lecchip/.style={lecchipbase}}%
  \ifnum#1=1 \tikzset{lecparta/.style={lecpartcur}}\fi
  \ifnum#1=2 \tikzset{lecpartb/.style={lecpartcur}}\fi
  \ifnum#1=3 \tikzset{lecpartc/.style={lecpartcur}}\fi
  \ifnum#1=4 \tikzset{lecpartd/.style={lecpartcur}}\fi
  \ifnum#1=5 \tikzset{lecparte/.style={lecpartcur}}\fi
  \lecmapchip{chipTone}{1}\lecmapchip{chipTtwoa}{2a}\lecmapchip{chipTtwob}{2b}%
  \lecmapchip{chipTthreea}{3a}\lecmapchip{chipTthreeb}{3b}\lecmapchip{chipTfour}{4}%
  \lecmapchip{chipTfivea}{5a}\lecmapchip{chipTfiveb}{5b}%
  \lecmapchip{chipTfivec}{5c}\lecmapchip{chipTfived}{5d}%
\fi
\begin{center}
\begin{tikzpicture}[
    part/.style={rectangle, rounded corners, draw=black!60, fill=VeryLightGray,
        minimum width=2.62cm, minimum height=0.72cm, text width=2.5cm,
        align=center, font=\scriptsize, text=DarkText},
    lecpartcur/.style={draw=RedTitle, very thick, fill=orange!20},
    lecpartdone/.style={draw=green!45!black, thick, fill=green!10},
    lecchipbase/.style={rectangle, rounded corners=1pt, draw=black!30, fill=white,
        inner xsep=1.5pt, inner ysep=1pt, font=\tiny, text=black!70},
    lecchipcur/.style={draw=RedTitle, fill=RedTitle!12, text=RedTitle, font=\tiny\bfseries},
    lecchipdone/.style={draw=green!45!black, fill=green!8},
    q/.style={font=\tiny\itshape, text width=2.7cm, align=center, text=black!75},
    twin/.style={font=\tiny, text width=2.7cm, align=center, text=blue!40!black},
    sat/.style={rectangle, rounded corners, draw=black!45, dashed, fill=white,
        text width=5.6cm, align=center, font=\tiny, text=black!70, minimum height=0.42cm},
    barlab/.style={font=\tiny\bfseries, text=black!60},
    arr/.style={-{Stealth[length=2mm]}, thick, gray!60}
]
    % --- five part boxes ---
    \node[part, lecparta] (p1) at (0,0)     {\textbf{9.1 SEE}};
    \node[part, lecpartb] (p2) at (2.85,0)  {\textbf{9.2 UNDER THE HOOD}};
    \node[part, lecpartc] (p3) at (5.7,0)   {\textbf{9.3 BUILD}};
    \node[part, lecpartd] (p4) at (8.55,0)  {\textbf{9.4 DESIGN \& VALIDATE}};
    \node[part, lecparte] (p5) at (11.4,0)  {\textbf{9.5 RUN \& GOVERN}};
    \draw[arr] (p1) -- (p2);
    \draw[arr] (p2) -- (p3);
    \draw[arr] (p3) -- (p4);
    \draw[arr] (p4) -- (p5);
    % --- sub-deck chips ---
    \node[lecchip, chipTone]    at (0,-0.6)     {T1};
    \node[lecchip, chipTtwoa]   at (2.55,-0.6)  {T2a};
    \node[lecchip, chipTtwob]   at (3.15,-0.6)  {T2b};
    \node[lecchip, chipTthreea] at (5.4,-0.6)   {T3a};
    \node[lecchip, chipTthreeb] at (6.0,-0.6)   {T3b};
    \node[lecchip, chipTfour]   at (8.55,-0.6)  {T4};
    \node[lecchip, chipTfivea]  at (10.5,-0.6)  {T5a};
    \node[lecchip, chipTfiveb]  at (11.1,-0.6)  {T5b};
    \node[lecchip, chipTfivec]  at (11.7,-0.6)  {T5c};
    \node[lecchip, chipTfived]  at (12.3,-0.6)  {T5d};
    % --- the student question each part answers ---
    \node[q] at (0,-1.08)     {what changed,\\ and why care?};
    \node[q] at (2.85,-1.08)  {how does it run?\\ how do I start?};
    \node[q] at (5.7,-1.08)   {how do I build one,\\ and call it?};
    \node[q] at (8.55,-1.08)  {can I design one\\ and prove it works?};
    \node[q] at (11.4,-1.08)  {run a project, catch\\ mistakes, stay safe};
    % --- its economic twin ---
    \node[twin] at (0,-1.62)     {generation got cheap;\\ verification is scarce};
    \node[twin] at (2.85,-1.62)  {what am I paying for?\\ tokens, scarce context};
    \node[twin] at (5.7,-1.62)   {dividing the labour:\\ specialists, boundaries};
    \node[twin] at (8.55,-1.62)  {write the contract,\\ audit the work};
    \node[twin] at (11.4,-1.62)  {hidden action: gates,\\ monitoring, priors};
    % --- "you are here" marker above the current part ---
    \ifnum#1=1 \node[font=\scriptsize\bfseries, text=RedTitle] at (0,0.62)
        {you are here\ $\blacktriangledown$};\fi
    \ifnum#1=2 \node[font=\scriptsize\bfseries, text=RedTitle] at (2.85,0.62)
        {you are here\ $\blacktriangledown$};\fi
    \ifnum#1=3 \node[font=\scriptsize\bfseries, text=RedTitle] at (5.7,0.62)
        {you are here\ $\blacktriangledown$};\fi
    \ifnum#1=4 \node[font=\scriptsize\bfseries, text=RedTitle] at (8.55,0.62)
        {you are here\ $\blacktriangledown$};\fi
    \ifnum#1=5 \node[font=\scriptsize\bfseries, text=RedTitle] at (11.4,0.62)
        {you are here\ $\blacktriangledown$};\fi
    \ifnum#1=0 \node[font=\scriptsize\bfseries, text=green!45!black] at (5.7,0.62)
        {all five parts complete};\fi
    % --- bottom band: the three questions of the lecture ---
    \fill[blue!10]   (-1.3,-2.37) rectangle (4.15,-2.07);
    \fill[green!10]  (4.4,-2.37)  rectangle (9.85,-2.07);
    \fill[orange!15] (10.1,-2.37) rectangle (12.7,-2.07);
    \node[barlab] at (1.425,-2.22)  {HOW IT WORKS};
    \node[barlab] at (7.125,-2.22)  {HOW TO BUILD \& USE IT};
    \node[barlab] at (11.4,-2.22)   {YOUR ROLE};
    % --- dashed satellites ---
    \node[sat] at (1.9,-2.8)
        {\textbf{Appendix TA} --- the dated landscape and setup details, read on demand};
    \node[sat] at (9.9,-2.8)
        {\textbf{Lab} Parts A--B, Design Lab scenarios $\to$ \textbf{Lec10} cases (same morning)};
\end{tikzpicture}
\end{center}
}


\title[AI for Econ Research]{AI for Economic Research: Dynamic Models, Language, and Agents\\
Lecture 9.3b: Invoking Agents --- One, Three, Ten Thousand}

\begin{document}

\makebeamertitle

% ==== Phase B · T2c-F01 · TIER: READ (map frame done in Phase A)
\begin{frame}{Lecture 9 in One Map: Invoking Agents}
\textbf{T3a left three job descriptions in a folder. This deck answers what you type to make one run --- and where that leads.}

\LecNineMap{3}{b}

\begin{takeawaybox}
\footnotesize\textbf{The claim this deck defends.} The number of agents you run is not a capability decision. It is a \emph{verification} decision --- you may only buy scale when a wrong answer is cheap to detect.
\end{takeawaybox}

{\scriptsize \textit{Reading convention (Appendix TA): \textbf{[solid]} $=$ primary or citable; \textbf{[hype]} $=$ vendor self-report.}}
\end{frame}


%=============================================================================
\section{Many Agents}
%===========================================================

\sectiondivider{Many Agents}{Creating a set of them --- and what actually limits the number}

% ---- T2b-F22 · KEEP · TIER: READ · from T2b:873
\begin{frame}{Three Dials Make an Agent}
\textbf{To create a \emph{set} of agents you do not create anything. You vary three dials and keep three separate lists.}

\vspace{0.12cm}

\begin{center}
\begin{tikzpicture}[
    dial/.style={rectangle, draw=RedTitle!60, fill=blue!5, rounded corners,
                 minimum width=2.5cm, minimum height=0.62cm, font=\scriptsize, align=center},
    ag/.style={rectangle, draw=black!60, fill=VeryLightGray, rounded corners,
               minimum width=2.75cm, minimum height=1.35cm, text width=2.6cm,
               align=center, font=\tiny},
    arr/.style={-{Stealth[length=1.8mm]}, thick, gray!65}
]
    \node[dial] (d1) at (0,1.15) {\textbf{system} --- identity};
    \node[dial] (d2) at (0,0.35) {\textbf{tools} --- capability};
    \node[dial] (d3) at (0,-0.45) {\textbf{messages} --- memory};

    \node[ag] (A) at (4.3, 1.05) {\textbf{Analyst}\\ stats tools\\ own list};
    \node[ag] (B) at (4.3,-0.55) {\textbf{Reviewer}\\ read-only tools\\ own list};
    \node[ag] (C) at (7.6, 0.25) {\textbf{Verifier}\\ test-runner tool\\ own list};

    \draw[arr] (1.35,0.35) -- (2.9,0.9);
    \draw[arr] (1.35,0.35) -- (2.9,-0.4);
    \draw[arr] (5.7,0.5) -- (6.2,0.35);

    \node[font=\scriptsize, align=left, text width=3.1cm] at (10.6,0.25)
        {Three agents $=$\\ three \emph{configurations}\\ and three \emph{lists}.\\[4pt]
         No processes, no\\ registry, no server\\ objects.};
\end{tikzpicture}
\end{center}

\vspace{0.05cm}

\begin{takeawaybox}
\footnotesize\textbf{The economist's reading.} You are not hiring staff; you are writing three job descriptions and keeping three notebooks. The ``firm'' exists only in your code. What is scarce is not agents --- it is tokens, context, and your attention.
\end{takeawaybox}
\end{frame}

% ---- T2b-F23 · EDIT · TIER: READ · from T2b:912
\begin{frame}{How Many at Once? Three Constraints, Not Infinity}
\textbf{``You can instantiate an unlimited number of agents'' is true in the same sense that you can write an unlimited number of cheques.}

\vspace{0.1cm}

\begin{center}
\begin{tikzpicture}[font=\scriptsize]
    \fill[green!8] (0,0) rectangle (3.4,1.55);
    \draw[-{Stealth[length=2mm]}, thick, black!70] (0,0) -- (11.2,0)
        node[right, font=\tiny] {agents in flight};
    \foreach \x/\l in {0/0, 1.7/2, 3.4/4, 5.1/6, 6.8/8, 8.5/10, 10.2/12}
        \draw[black!45] (\x,0.06) -- (\x,-0.06) node[below, font=\tiny] {\l};

    \draw[dashed, thick, red!65!black] (3.4,-0.2) -- (3.4,1.9);
    \node[align=center, font=\tiny, red!55!black, anchor=south] at (3.4,1.92)
        {\textbf{tokens/min}\\ rate limit};

    \draw[dashed, thick, orange!75!black] (6.0,-0.2) -- (6.0,1.45);
    \node[align=center, font=\tiny, orange!65!black, anchor=south] at (6.2,1.47)
        {\textbf{context}\\ each agent's list grows};

    \draw[dashed, thick, RedTitle!75] (8.7,-0.2) -- (8.7,1.0);
    \node[align=center, font=\tiny, RedTitle, anchor=south] at (8.9,1.02)
        {\textbf{coordination}\\ merge + verify cost};

    \node[font=\tiny, align=center, green!35!black] at (1.7,0.72)
        {feasible\\ and useful};
\end{tikzpicture}
\end{center}

\vspace{0.05cm}

\begin{itemize}\footnotesize
    \item \textbf{Rate limits} are the hard ceiling: quota is tokens-per-minute, not agents. Ten agents each re-sending a long list exhaust it faster than two.
    \item \textbf{Context} binds per agent: every worker carries its own growing list (previous section).
    \item \textbf{Coordination} is the one economists should expect: \emph{you} must read, reconcile, and verify $n$ outputs. That cost rises with $n$ while the value of the marginal agent falls.
\end{itemize}

\vspace{0.1cm}
\footnotesize\textbf{The binding constraint is usually the third one.} T1 priced this out: multi-agent bought a large accuracy gain at roughly $15\times$ the token cost. Reach for more agents when the task is genuinely wide, not to feel parallel.
\end{frame}

% ---- T2b-F24 · EDIT · TIER: CORE · from T2b:954
\begin{frame}{Isolated Context Is the Point}
\textbf{The reason to spawn a sub-agent is not extra brains --- the model is the same. It is a second, clean list.}

\vspace{0.02cm}

\begin{center}
\begin{tikzpicture}[font=\scriptsize, yscale=0.80]
    \draw[draw=black!60, fill=VeryLightGray, rounded corners] (0,0) rectangle (3.3,2.6);
    \fill[RedTitle!22] (0.06,0.06) rectangle (3.24,2.05);
    \node[font=\scriptsize\bfseries, align=center] at (1.65,2.88) {Orchestrator};
    \node[font=\tiny, align=center] at (1.65,1.05)
        {plan, six file reads,\\ three tool results,\\ two dead ends,\\ your last nine prompts};
    \node[font=\tiny, align=center] at (1.65,2.32) {\emph{78\% full}};

    \draw[draw=black!60, fill=VeryLightGray, rounded corners] (5.2,0) rectangle (7.7,2.6);
    \fill[green!45!black!30] (5.26,0.06) rectangle (7.64,0.55);
    \node[font=\scriptsize\bfseries, align=center] at (6.45,2.88) {Worker A};
    \node[font=\tiny, align=center] at (6.45,0.3) {one task, one file};
    \node[font=\tiny, align=center, gray!70] at (6.45,1.5) {\emph{clean}};

    \draw[draw=black!60, fill=VeryLightGray, rounded corners] (9.0,0) rectangle (11.5,2.6);
    \fill[green!45!black!30] (9.06,0.06) rectangle (11.44,0.5);
    \node[font=\scriptsize\bfseries, align=center] at (10.25,2.88) {Worker B};
    \node[font=\tiny, align=center] at (10.25,0.28) {one task, one file};
    \node[font=\tiny, align=center, gray!70] at (10.25,1.5) {\emph{clean}};

    \draw[-{Stealth[length=2mm]}, thick, gray!70] (3.4,1.9) -- (5.1,1.9)
        node[midway, above, font=\tiny] {task only};
    \draw[-{Stealth[length=2mm]}, thick, gray!70] (3.4,1.9) to[bend left=8] (8.9,2.0);

    \draw[-{Stealth[length=2mm]}, thick, green!45!black] (6.45,-0.35) -- (1.65,-0.35)
        node[midway, below, font=\tiny] {\emph{findings only} --- a short summary, not the transcript};
    \draw[green!45!black, thick] (6.45,0) -- (6.45,-0.35);
    \draw[green!45!black, thick] (10.25,0) -- (10.25,-0.35) -- (6.45,-0.35);
    \draw[green!45!black, thick] (1.65,-0.35) -- (1.65,0);
\end{tikzpicture}
\end{center}

\vspace{0.22cm}

\begin{columns}[T]
\begin{column}{0.55\textwidth}
\begin{itemize}\footnotesize
    \item The worker never sees your dead ends, so cannot be misled by them
    \item The orchestrator's list grows by a summary, not by all the worker read
    \item That asymmetry is the economic case for sub-agents
\end{itemize}
\end{column}
\begin{column}{0.43\textwidth}
\begin{cautionbox}
\footnotesize\textbf{Do not over-read this.} A clean window is not an independent mind. T5c shows workers inherit your framing --- fresh context buys \emph{focus}, not \emph{independence}.
\end{cautionbox}
\end{column}
\end{columns}
\end{frame}

% ---- T2b-F25 · EDIT · TIER: READ · from T2b:1010
\begin{frame}{Three Wirings, Drawn as Data Flow}
\textbf{Strip the vocabulary away and multi-agent designs differ in exactly one respect: who writes into whose list.}

\vspace{0.15cm}

\begin{center}
\begin{tikzpicture}[
    n/.style={circle, draw=black!65, fill=VeryLightGray, minimum size=0.66cm, font=\tiny},
    bb/.style={rectangle, draw=RedTitle!70, fill=blue!8, rounded corners,
               minimum width=2.2cm, minimum height=0.5cm, font=\tiny},
    arr/.style={-{Stealth[length=1.7mm]}, thick, gray!70},
    hd/.style={font=\scriptsize\bfseries}
]
    % Sequential
    \node[hd] at (1.5,1.5) {Sequential};
    \node[n] (s1) at (0.3,0.4) {A};
    \node[n] (s2) at (1.5,0.4) {B};
    \node[n] (s3) at (2.7,0.4) {C};
    \draw[arr] (s1) -- (s2);
    \draw[arr] (s2) -- (s3);
    \node[font=\tiny, align=center, gray!75] at (1.5,-0.45) {each list seeded\\ by the last output};

    % Hierarchical
    \node[hd] at (5.9,1.5) {Hierarchical};
    \node[n] (h0) at (5.9,0.85) {O};
    \node[n] (h1) at (4.9,-0.1) {A};
    \node[n] (h2) at (5.9,-0.1) {B};
    \node[n] (h3) at (6.9,-0.1) {C};
    \draw[arr] (h0) -- (h1);
    \draw[arr] (h0) -- (h2);
    \draw[arr] (h0) -- (h3);
    \node[font=\tiny, align=center, gray!75] at (5.9,-0.8) {workers report up;\\ only O sees all};

    % Blackboard
    \node[hd] at (10.0,1.5) {Blackboard};
    \node[bb] (bbn) at (10.0,0.85) {shared state};
    \node[n] (b1) at (9.0,-0.15) {A};
    \node[n] (b2) at (10.0,-0.15) {B};
    \node[n] (b3) at (11.0,-0.15) {C};
    \draw[arr, <->] (bbn) -- (b1);
    \draw[arr, <->] (bbn) -- (b2);
    \draw[arr, <->] (bbn) -- (b3);
    \node[font=\tiny, align=center, gray!75] at (10.0,-0.85) {all read and write\\ one store};
\end{tikzpicture}
\end{center}

\vspace{0.35cm}

\begin{itemize}\footnotesize
    \item \textbf{Sequential} is a relay: cheap, easy to debug, no parallel gain, and an early error propagates
    \item \textbf{Hierarchical} is T1's orchestrator--worker and what T5b's RA team is --- the pattern that dominates in practice
    \item \textbf{Blackboard} decouples agents from each other but needs write discipline; in a research project the ``blackboard'' is usually just \emph{files on disk} --- exactly the file-driven workflow of T5b
\end{itemize}
\end{frame}


%=============================================================================
\section{Who Decides to Dispatch?}
%===========================================================

\sectiondivider{Who Decides to Dispatch?}{The router myth, and what the transcript actually shows}

% ---- T2b-F26 · EDIT · TIER: READ · from T2b:1071
\begin{frame}{The Router Myth}
\textbf{The usual explanation: a fast ``router'' model classifies your prompt, looks up specialists in an agent registry, and wakes them. That is not what happens.}

\vspace{0.12cm}

\begin{center}
\begin{tikzpicture}[
    bx/.style={rectangle, draw=black!60, fill=VeryLightGray, rounded corners,
               minimum width=2.0cm, minimum height=0.6cm, font=\tiny, align=center},
    gh/.style={rectangle, draw=red!55!black, fill=red!5, rounded corners,
               minimum width=2.0cm, minimum height=0.6cm, font=\tiny, align=center},
    ok/.style={rectangle, draw=green!45!black, fill=green!6, rounded corners,
               minimum width=2.0cm, minimum height=0.6cm, font=\tiny, align=center},
    arr/.style={-{Stealth[length=1.7mm]}, thick, gray!70},
    hd/.style={font=\scriptsize\bfseries}
]
    % LEFT: the myth
    \node[hd, red!55!black] at (2.4,2.35) {What people describe};
    \node[bx] (u1) at (0.0,1.5) {your prompt};
    \node[gh] (rt) at (2.4,1.5) {\textbf{router LLM}};
    \node[gh] (rg) at (2.4,0.5) {agent registry};
    \node[bx] (w1) at (4.8,1.9) {specialist 1};
    \node[bx] (w2) at (4.8,1.05) {specialist 2};
    \draw[arr] (u1) -- (rt);
    \draw[arr] (rt) -- (rg);
    \draw[arr] (rt) -- (w1);
    \draw[arr] (rt) -- (w2);
    \draw[red!60!black, very thick, opacity=0.55] (-1.0,-0.15) -- (5.9,2.6);
    \draw[red!60!black, very thick, opacity=0.55] (-1.0,2.6) -- (5.9,-0.15);

    % RIGHT: reality
    \node[hd, green!35!black] at (9.6,2.35) {What actually happens};
    \node[bx] (u2) at (7.3,1.5) {your prompt};
    \node[ok] (md) at (9.7,1.5) {\textbf{same model,}\\ \textbf{same context}};
    \node[ok] (tc) at (9.7,0.5) {emits a tool call\\ named \texttt{Agent}};
    \node[bx] (sw) at (12.1,0.5) {sub-agent\\ runs};
    \draw[arr] (u2) -- (md);
    \draw[arr] (md) -- (tc);
    \draw[arr] (tc) -- (sw);
\end{tikzpicture}
\end{center}

\vspace{0.05cm}

\begin{itemize}\footnotesize
    \item \textbf{No second model, no classification step.} Dispatch is one more \texttt{tool\_use} block from T2a's loop --- the tool just happens to start another loop
    \item The ``registry'' is a folder of Markdown files \emph{you wrote}; the model picks from it by reading their \texttt{description} fields, exactly as it picks any tool
    \item T2a said this for one agent: \emph{``the same transformer that writes code decides to call tools.''} This is that fact one level up
\end{itemize}

{\scriptsize \textit{Guard: TA's ``think-budget routers'' route one query between fast and deep \emph{modes of one model}, not between agents.}}
\end{frame}

% ==== Phase B · T2b-F27 · TIER: READ (observation provenance kept)
\begin{frame}{Dispatch Is a Tool Call --- Evidence You Can See}
\textbf{You do not have to take this on faith. Four checks, all runnable in the lab that follows.}

\vspace{0.12cm}

\begin{center}
\footnotesize
\renewcommand{\arraystretch}{1.22}
\begin{tabular}{|p{0.4cm}|p{4.3cm}|p{5.9cm}|}
\hline
& \textbf{Do this} & \textbf{What it proves} \\ \hline
1 & Run \texttt{/agents} in a session & The list is \emph{your} \texttt{.claude/agents/} folder. Nothing is registered with a service. \\ \hline
2 & Ask something a specialist fits; watch the transcript & The dispatch appears as a named tool call in the same stream as \texttt{Read} and \texttt{Bash}. \\ \hline
3 & Compare context before and after & It grows by a short \emph{summary}, not by the worker's transcript. \\ \hline
4 & \textbf{Narrow one agent's \texttt{description}, re-ask} & Dispatch stops. Selection is driven by prose the model reads. \\ \hline
\end{tabular}
\end{center}

\vspace{0.12cm}

\begin{takeawaybox}
\footnotesize\textbf{Check 4 is the decisive one.} If a separate router with a registry were making the decision, editing one English sentence in a Markdown file could not switch the behaviour off. It can --- because that sentence \emph{is} the mechanism.
\end{takeawaybox}

\vspace{0.06cm}
{\scriptsize \textit{Real trace, check 2 (Claude Code 2.1.261): \texttt{Read} $\to$ \texttt{Agent code-reviewer} $\to$ \texttt{Agent domain-reviewer} --- the dispatch sits in the same list as \texttt{Read}. \textbf{[solid]}}}
\end{frame}

% ---- T2b-F28 · EDIT · TIER: READ · from T2b:1152
\begin{frame}{Auto-Dispatch vs.\ Hand-Wired: Spontaneous Order or Central Planning}
\textbf{Once you know dispatch is a model decision, the design choice becomes familiar: leave allocation to the agent, or plan it yourself.}

\vspace{0.05cm}

\begin{center}
\begin{tikzpicture}[font=\scriptsize, yscale=0.86]
    \draw[-{Stealth[length=2mm]}, thick, black!65] (0,0) -- (11.4,0);
    \node[font=\tiny, anchor=north] at (0.9,-0.1) {you specify everything};
    \node[font=\tiny, anchor=north] at (10.3,-0.1) {the model decides};

    \node[rectangle, draw=RedTitle!65, fill=blue!6, rounded corners, minimum width=3.2cm,
          minimum height=1.2cm, text width=3.0cm, align=center, font=\tiny] (L) at (1.9,1.1)
        {\textbf{Hand-wired}\\ fixed pipeline you author\\ \emph{low variance, high fixed cost}};
    \node[rectangle, draw=black!55, fill=VeryLightGray, rounded corners, minimum width=3.2cm,
          minimum height=1.2cm, text width=3.0cm, align=center, font=\tiny] (M) at (5.7,1.1)
        {\textbf{Where research sits}\\ you author the specialists,\\ the model chooses when};
    \node[rectangle, draw=green!45!black, fill=green!5, rounded corners, minimum width=3.2cm,
          minimum height=1.2cm, text width=3.0cm, align=center, font=\tiny] (R) at (9.5,1.1)
        {\textbf{Auto-dispatch}\\ agent picks its own helpers\\ \emph{adapts, variable cost}};
    \draw[gray!60, dashed] (1.9,0.48) -- (1.9,0.04);
    \draw[gray!60, dashed] (5.7,0.48) -- (5.7,0.04);
    \draw[gray!60, dashed] (9.5,0.48) -- (9.5,0.04);
\end{tikzpicture}
\end{center}

\vspace{0.1cm}

\begin{columns}[T]
\begin{column}{0.49\textwidth}
\footnotesize\textbf{Central planning (hand-wired)}
\begin{itemize}\footnotesize
    \item Reproducible: same path every run
    \item No planner tokens spent
    \item Brittle off the anticipated path
    \item High authoring cost; it decays
\end{itemize}
\end{column}
\begin{column}{0.49\textwidth}
\footnotesize\textbf{Spontaneous order (auto-dispatch)}
\begin{itemize}\footnotesize
    \item Handles tasks you did not foresee
    \item Zero setup at the point of use
    \item \textbf{Principal--agent problem:} it may delegate the wrong thing, and you pay
    \item Variable, harder-to-forecast cost
\end{itemize}
\end{column}
\end{columns}

\vspace{0.1cm}
\footnotesize\textbf{The practical middle.} Author the specialists once, let the model choose \emph{when} to call them, and keep a validation gate so a bad delegation is caught rather than trusted --- T5b's quality gates, doing an economist's job.
\end{frame}


%=============================================================================
\section{Three Rungs}
%===========================================================

\sectiondivider{Three Rungs}{The only three ways to invoke an agent --- and who chooses in each}

% ---- T2c-F02 · EDIT · TIER: CORE · from T2c:72
\begin{frame}{The Question T3a Left Open}
\textbf{Three job descriptions sit in a folder. Nothing yet runs them. There are exactly three ways to start one, and they differ in one thing: who chooses.}

\vspace{0.14cm}

\begin{center}
\begin{tikzpicture}[font=\scriptsize, yscale=0.9,
    rung/.style={rectangle, draw=black!60, fill=VeryLightGray, rounded corners,
                 minimum width=3.5cm, minimum height=1.15cm, text width=3.3cm, align=center},
    hi/.style={rung, draw=RedTitle!70, fill=blue!6},
    cpt/.style={font=\tiny, align=center, gray!75, text width=3.4cm}]
    \node[rung] (r1) at (0,0)      {\textbf{1. Implicit}\\ you describe the task};
    \node[rung] (r2) at (4.0,0.95) {\textbf{2. Explicit}\\ you name the agent};
    \node[hi]   (r3) at (8.0,1.9)  {\textbf{3. Scripted}\\ you loop over processes};
    \node[cpt] at (0,-1.05)   {\emph{the model} chooses, by reading \texttt{description} fields\\ \textbf{reaches 1--3}};
    \node[cpt] at (4.0,-0.10)  {\emph{you} choose\\ \textbf{reaches 1--5}};
    \node[cpt] at (8.0,0.85)   {nobody chooses --- the loop runs all $N$\\ \textbf{reaches $10^4$}};
    \draw[-{Stealth[length=1.8mm]}, thick, gray!65] (1.8,0.3) -- (2.2,0.65);
    \draw[-{Stealth[length=1.8mm]}, thick, gray!65] (5.8,1.25) -- (6.2,1.6);
\end{tikzpicture}
\end{center}

\vspace{0.1cm}

\begin{takeawaybox}
\footnotesize\textbf{Only the third rung scales, because only it removes you from the loop.} Rungs 1 and 2 cost a turn of your attention each --- the third ceiling above, the one that binds. Rung 3 trades that for a prior problem: what checks the output.
\end{takeawaybox}

\vspace{0.06cm}
{\footnotesize\textbf{Answered next:} the exhibit the router section argued for but never showed.}
\end{frame}

% ==== Phase B · T2c-F03 · TIER: LAB (harness line)
\begin{frame}[fragile]{Rung 1 --- You Describe It, the Model Chooses}
\textbf{The router section proved dispatch is an ordinary tool call, by argument and by experiment. It never showed you one. Here it is.}

\vspace{0.08cm}

\begin{columns}[T]
\begin{column}{0.55\textwidth}
\begin{lstlisting}[style=pythonstyle,language={},basicstyle=\ttfamily\fontsize{5.2}{6.1}\selectfont,numbers=none]
you  >  Check the derivation in src/aiyagari.py
        against docs/model.md.

# what the model emits -- one block in the same
# list as Read and Bash, nothing special about it
{ "type": "tool_use",
  "name": "Agent",
  "input": {
    "subagent_type": "math-agent",
    "description": "check the derivation",
    "prompt": "Review src/aiyagari.py against
               docs/model.md. The calibration is
               annual. Report only findings you
               can tie to a line number." } }

# ... the sub-agent's whole session happens here ...

{ "type": "tool_result",
  "content": "line 34: utility() has no domain
              guard. With SIGMA=2, u(c)=-1/c is
              positive for c<0 ..." }
\end{lstlisting}
\end{column}
\begin{column}{0.43\textwidth}
\begin{itemize}\footnotesize
  \item \textbf{You typed no agent name.} The model read the \texttt{description} fields in your folder and chose --- the same way it chooses \texttt{Read}
  \item \textbf{\texttt{prompt} is the brief}, T3a's step 4, written by the model this time rather than by you
  \item \textbf{One \texttt{tool\_result} comes back.} Everything the sub-agent read is gone --- T3a's step 5
\end{itemize}

\vspace{0.1cm}
{\scriptsize \textit{The same thing in a real transcript, from the lab: \texttt{Bash} $\to$ \texttt{Agent code-reviewer} $\to$ \texttt{Read} --- \texttt{Agent} sits in the same list as the others. That is the whole claim. \textbf{[solid]}}}
\end{column}
\end{columns}

\vspace{0.05cm}
{\footnotesize\textbf{Answered next:} naming it yourself.} {\tiny\textit{Elsewhere: \texttt{.gemini/agents/} (trusted folders), Codex's \texttt{spawn\_agent}; transcripts not yet verified.}}
\end{frame}

% ==== Phase B · T2c-F04 · TIER: READ (harness line)
\begin{frame}[fragile]{Rung 2 --- You Name It}
\textbf{Naming the agent moves the choice from the model to you. Three surfaces do it, and this course already uses two of them without saying so.}

\vspace{0.08cm}

\begin{columns}[T]
\begin{column}{0.56\textwidth}
\begin{lstlisting}[style=pythonstyle,language={},basicstyle=\ttfamily\fontsize{5.2}{6.1}\selectfont,numbers=none]
# (a) in prose -- the idiom the OLG demo already uses
you >  Send the pseudo-code to the domain-reviewer
       sub-agent; require it to trace one cohort's
       residual end to end before I write code.

# (b) inspect what is actually in scope
/agents

# (c) pin a whole session to one specialist
claude --agent math-agent

# (d) define one inline, with no file at all
claude --agents '{"claim-checker":
  {"description":"Find the source of every number.",
   "prompt":"You verify figures against the repo."}}'
\end{lstlisting}

\vspace{0.04cm}
{\scriptsize \textit{\texttt{/agents} is the one command T2b's twelve do not list; T2b owns the rest of the surface.}}
\end{column}
\begin{column}{0.42\textwidth}
\begin{cautionbox}
\footnotesize\textbf{Naming it skips the only test you had.} Rung 1 makes the model choose by reading your \texttt{description}. That choice \emph{is} the test that the description is right --- T3a's step 7 is built on it.

\vspace{0.1cm}
Name the agent every time and a badly worded \texttt{description} never fails, so you never find out. Use rung 2 when you know which specialist you want; use rung 1 when you are still debugging the English.
\end{cautionbox}
\end{column}
\end{columns}

\vspace{0.06cm}
{\footnotesize\textbf{Answered next:} the rung that removes you, and reaches four digits.} {\tiny\textit{\texttt{/agents}, \texttt{-{}-agent} and \texttt{-{}-agents} are Claude Code's surfaces; elsewhere, name the agent in prose or use rung 3.}}
\end{frame}

% ==== Phase B · T2c-F05 · TIER: READ (headless equivalents)
\begin{frame}[fragile]{Rung 3 --- You Script It}
\textbf{Rungs 1 and 2 spend a turn of your attention per dispatch. Rung 3 spends a process instead --- and it is the only rung that reaches four digits.}

\vspace{0.08cm}

\begin{columns}[T]
\begin{column}{0.56\textwidth}
\begin{lstlisting}[style=pythonstyle,language={},basicstyle=\ttfamily\fontsize{5.2}{6.1}\selectfont,numbers=none]
# one process per unit of work, fenced read-only
for f in source/Lec09_Agentic_AI/Lec09_T*.tex; do
    claude -p "$(cat briefs/claim_check.md) $f" \
           --output-format json \
           --allowedTools Read Grep Glob \
           --permission-mode dontAsk \
      > "out/$(basename "$f" .tex).json" &
done
wait
# or hand them to the background control plane
claude --bg "audit T3b for unsourced numbers"
claude agents          # list them
claude logs <id>       # read one
claude stop <id>       # kill one
\end{lstlisting}
\end{column}
\begin{column}{0.42\textwidth}
\begin{center}
\footnotesize
\renewcommand{\arraystretch}{1.2}
\begin{tabular}{|p{1.5cm}|p{1.5cm}|p{2.0cm}|}
\hline
\textbf{Rung} & \textbf{Costs you} & \textbf{Binds at} \\ \hline
1 implicit & a turn, each time & your attention \\ \hline
2 explicit & a turn, each time & your attention \\ \hline
3 scripted & one brief, once & tokens per minute \\ \hline
\end{tabular}
\end{center}

\vspace{0.1cm}

\begin{takeawaybox}
\footnotesize\textbf{The brief becomes a file.} At rung 3 you cannot retype it, so it must live on disk --- T3a's step 8, arriving by necessity rather than discipline.
\end{takeawaybox}
\end{column}
\end{columns}

\vspace{0.05cm}
{\tiny \textit{Checked on Claude Code 2.1.269; elsewhere \texttt{gemini -p "\ldots" -o json}, \texttt{codex exec "\ldots" -{}-json} (help, 12 Sep 2026). No turn-cap flag: budget via model, allowlist, brief.}}
\end{frame}

% ---- T2c-F14 · EDIT · TIER: CORE · from T2c:478
\begin{frame}{Choosing a Rung Is One Question}
\textbf{Not ``how many agents can I run'' --- The first section of this deck already answered that. The question is what checks the output, and how cheaply.}

\vspace{0.14cm}

\begin{center}
\footnotesize
\renewcommand{\arraystretch}{1.25}
\begin{tabular}{|p{4.3cm}|p{2.2cm}|p{4.2cm}|}
\hline
\textbf{What can check the output} & \textbf{Rung} & \textbf{Worked example} \\ \hline
only you, by reading & 1--2, $N$ small & the three background specialists, judged by a $3 \times 3$ detection matrix (T3a) \\ \hline
a cheap partial check --- tests, a rubric, a schema --- plus a tier for the hard cases & 3, $N$ in the thousands & T6b: cached rubric, Sonnet adjudication, 11{,}171 labels \\ \hline
a sound mechanical check & 3, $N$ until the quota binds & Lean, a compiler, a one-command replication script \\ \hline
\end{tabular}
\end{center}

\vspace{0.14cm}

\begin{center}
\textbf{$N$ is a verification choice, not a capability choice.}\\
{\footnotesize\textit{The ceilings tell you where the limit is. The verifier tells you whether you want to go near it.}}
\end{center}
\end{frame}


%=============================================================================
\section{Four-Digit Runs}
%===========================================================

\sectiondivider{Four-Digit Runs}{What a four-digit agent run consists of --- and what made it bankable}

% ==== Phase B · T2c-F06 · TIER: READ (dated sidebar)
\begin{frame}{A Dated Sidebar: What About Ten Thousand Agents?}
\textbf{On 28 August 2026 OpenAI pointed roughly ten thousand agents at Navier--Stokes and had a Lean-checked proof 88 hours later. The number is the least interesting part.}

\vspace{0.12cm}

\begin{center}
\footnotesize
\renewcommand{\arraystretch}{1.22}
\begin{tabular}{|p{2.9cm}|p{8.0cm}|}
\hline
Warm-up (Euler) & $\sim$100 agents, $\sim$50 hours \\ \hline
The run (Navier--Stokes) & $\sim$10{,}000 agents from 28 Aug 2026, done within \textbf{88 hours} \\ \hline
Formalisation & a further \textbf{17 hours} to check the proof in \textbf{Lean} \\ \hline
Traffic & $\sim$5 million messages \emph{between agents}; $\sim$300 billion output tokens \\ \hline
Compute & ``several million dollars'' (Bubeck, OpenAI) \\ \hline
What was shown & finite-time singularity from smooth data --- the \emph{negative} resolution \\ \hline
\end{tabular}
\end{center}

\vspace{0.1cm}

\begin{cautionbox}
\footnotesize\textbf{Read the ledger, not the headline.} Every figure is the vendor's own account of an unreleased model, and the result is disputed two frames on. \textbf{[hype]} What is worth your attention is the \emph{architecture} --- reported consistently, and the part that transfers.
\end{cautionbox}

\vspace{0.05cm}
{\scriptsize \textit{As of 12 Sep 2026 (announced 8 Sep; Quanta, Semafor, CNN) --- re-check before teaching. \textbf{[solid]} that it was claimed and is contested.}}
\end{frame}

% ---- T2c-F07 · EDIT · TIER: READ · from T2c:277
\begin{frame}{Not Ten Thousand Geniuses: Groups, Framings, and a Human}
\textbf{Not ten thousand independent attempts. Agents sat in groups that talked internally, each group got a \emph{different framing of the same problem}, and a human moved fragments between groups.}

\vspace{0.1cm}

\begin{center}
\begin{tikzpicture}[font=\tiny, yscale=0.82,
    grp/.style={rectangle, draw=black!60, fill=VeryLightGray, rounded corners,
                minimum width=2.5cm, minimum height=0.92cm, text width=2.35cm, align=center},
    par/.style={rectangle, draw=green!45!black, fill=green!6, rounded corners,
                minimum width=2.2cm, minimum height=0.8cm, text width=2.05cm, align=center},
    hum/.style={rectangle, draw=RedTitle!70, fill=blue!7, rounded corners,
                minimum width=2.4cm, minimum height=0.9cm, text width=2.25cm, align=center},
    arr/.style={-{Stealth[length=1.7mm]}, thick, gray!70},
    back/.style={-{Stealth[length=1.7mm]}, thick, RedTitle!65, dashed}]
    \node[grp] (g1) at (0,1.35)  {\textbf{group A}\\ framing 1};
    \node[grp] (g2) at (0,0.15)  {\textbf{group B}\\ framing 2};
    \node[grp] (g3) at (0,-1.05) {\textbf{group C}\\ framing 3};
    \node[par] (p)  at (4.0,0.15) {partial results};
    \node[hum] (h)  at (7.7,0.15) {\textbf{a human} reads and picks};
    \foreach \i in {1,2,3} {\draw[arr] (g\i) -- (p);}
    \draw[arr] (p) -- (h);
    \draw[back] (h.south) -- (7.7,-2.05) -- (-1.95,-2.05) -- (-1.95,0.15) -- (g2.west);
    \node[font=\tiny, RedTitle!75] at (3.2,-2.3) {re-seeded as new prompts};
    \node[font=\tiny, align=left, text width=2.6cm, gray!80] at (10.9,0.15)
        {within a group: agents talk.\\[3pt] across groups: only what a human carries.};
\end{tikzpicture}
\end{center}

\vspace{0.08cm}

\begin{takeawaybox}
\footnotesize\textbf{What varied was dial one.} Different framings of one problem \emph{are} different \texttt{system} prompts --- T3a's step 2, run ten thousand times instead of three --- wired as the blackboard above with a person as the board. Sober reading: a thousand-agent run with better synthesis might land in the same place.
\end{takeawaybox}
\end{frame}

% ---- T2c-F08 · EDIT · TIER: CORE · from T2c:313
\begin{frame}{The Verifier Is the Whole Story}
\textbf{You may only spend money on scale when a wrong answer is cheap to detect. Lean made a wrong proof cost nothing. That, and not the ten thousand, is the part that transfers.}

\vspace{0.1cm}

\begin{center}
\footnotesize
\renewcommand{\arraystretch}{1.2}
\begin{tabular}{|p{4.6cm}|p{6.3cm}|}
\hline
\textbf{Reporting} best-of-$n$ as your score & the sin T4 names: $\tau$-bench pass-hat-$k$ falls from $\sim$90\% at $k{=}1$ to $\sim$57\% at $k{=}8$, so a best-of-$n$ headline flatters the model and hides the reliability curve \\ \hline
\textbf{Running} $n$ attempts behind a checker that can reject & legitimate --- because the \emph{checker} decides, not the score. The $9{,}999$ failures cost tokens and nothing else \\ \hline
\end{tabular}
\end{center}

\vspace{0.1cm}

\begin{cautionbox}
\footnotesize\textbf{Two things Lean does not buy.} A human must still certify that the statement checked in Lean is the theorem anyone meant to prove --- the formal check is sound about the object it was given, not about your intent. And the result is contested: Buckmaster alleges OpenAI muscled in on work he and Alp\"oge were finishing, called one AI-written paper ``AI slop'' \emph{despite} its formal verification, and OpenAI denies the charge. \textbf{[solid]} on the dispute.
\end{cautionbox}

\vspace{0.05cm}
{\footnotesize\textbf{Answered next:} you already own a run like this --- and it worked for exactly this reason.}
\end{frame}

% ---- T2c-F09 · KEEP · TIER: READ · from T2c:338
\begin{frame}{Your Own Four-Digit Run, and Why It Worked}
\textbf{Lec10's T6b labelled 11{,}171 of 13{,}225 tasks in a single five-hour window. Nobody called it a swarm, but it is the same machine --- and it had a verifier.}

\vspace{0.12cm}

\begin{center}
\footnotesize
\renewcommand{\arraystretch}{1.2}
\begin{tabular}{|p{3.1cm}|p{3.8cm}|p{3.8cm}|}
\hline
 & \textbf{Navier--Stokes} & \textbf{T6b, your labelling run} \\ \hline
The checker & Lean: sound, mechanical & a rubric, cached per task \\ \hline
The hard cases & re-framed and re-seeded by a human & a Sonnet \textbf{adjudication tier} for the 3{,}456 flagged \\ \hline
Not repeating work & --- & one cache file per task; re-runs skip \\ \hline
Human time & re-seeding between groups & $\sim$20 min to spec the rubric; $\sim$12 labels read by eye \\ \hline
\end{tabular}
\end{center}

\vspace{0.1cm}

\begin{takeawaybox}
\footnotesize\textbf{Rubric $+$ adjudication tier $+$ cache \emph{is} a verifier} --- weaker than Lean, cheap enough to run 13{,}000 times. \textbf{And note what neither buys:} you can re-run a proof attempt or a labelling task ten thousand times, but you cannot re-run 2008. A central bank gets one history, so fan-out never manufactures a counterfactual (Lec05).
\end{takeawaybox}

\vspace{0.05cm}
{\scriptsize \textit{T6b owns these numbers and the bulking/caching design; this frame owns only the mapping. \textbf{[solid]}}}
\end{frame}


%=============================================================================
\section{The Ledger}
%===========================================================

\sectiondivider{The Ledger}{What four sub-agents cost --- the bill, and the design work}

% ---- T2c-F10 · EDIT · TIER: READ · from T2c:372
\begin{frame}{The Ledger: What Four Sub-Agents Cost}
\textbf{The thirteen frames of T3a's eight-step section were not written by one session. Four agents did the reading. Here is the bill, from this project's own logs.}

\vspace{0.12cm}

\begin{center}
\footnotesize
\renewcommand{\arraystretch}{1.2}
\begin{tabular}{|p{1.3cm}|p{5.1cm}|p{1.7cm}|p{1.1cm}|p{1.1cm}|}
\hline
\textbf{Type} & \textbf{Brief} & \textbf{Tokens} & \textbf{Calls} & \textbf{Min} \\ \hline
Explore & map the target deck and its house style & 57{,}474 & 15 & 3.0 \\ \hline
Explore & build system, CJK support, TikZ idioms & 121{,}309 & 43 & 5.0 \\ \hline
Explore & course context, capstone, PDF tooling & 145{,}432 & 34 & 5.6 \\ \hline
Plan & draft the frame-by-frame outline & 124{,}361 & 30 & 10.3 \\ \hline
\multicolumn{2}{|r|}{\textbf{total}} & \textbf{448{,}576} & \textbf{122} & \textbf{23.8} \\ \hline
\end{tabular}
\end{center}

\vspace{0.12cm}

\begin{takeawaybox}
\footnotesize\textbf{448{,}576 tokens were read; four summaries came back} --- roughly twenty- to thirty-fold compression, and the orchestrating session never held the 122 file reads that produced them. That is the isolated-context argument above with a price on it. The three explorers also ran \emph{concurrently}: \textbf{5.6 minutes of wall clock against 13.5 sequential} (T5b owns the general speed claim).
\end{takeawaybox}

\vspace{0.05cm}
{\scriptsize \textit{Own session logs, 12 September 2026, Claude Code 2.1.269. \textbf{[solid]}}}
\end{frame}

% ---- T2c-F11 · EDIT · TIER: READ · from T2c:401
\begin{frame}{The Design Work Was the Four Briefs}
\textbf{Dispatching cost one line each. What decided whether the deck was any good was what each brief asked for --- T3a's step 4, done four times.}

\vspace{0.12cm}

\begin{center}
\footnotesize
\renewcommand{\arraystretch}{1.2}
\begin{tabular}{|p{3.0cm}|p{3.9cm}|p{3.7cm}|}
\hline
\textbf{Asked for} & \textbf{Came back} & \textbf{Changed} \\ \hline
sections, frames, verbatim style patterns & no deck shows an agent file & the section's premise, and the style contract \\ \hline
engine, CJK support, best existing figures & a Chinese guard exists in Lec01; zero \texttt{includegraphics}; two TikZ idioms & staying on \texttt{pdflatex}; both card figures \\ \hline
capstone, prerequisites, PDF tooling & prerequisites require neither Python nor ML; a converter pipeline exists on disk & claims 1 and 4; the pipeline frame \\ \hline
\end{tabular}
\end{center}

\vspace{0.12cm}

\begin{takeawaybox}
\footnotesize\textbf{None of the three asked a question whose answer was already known.} Each asked for an \emph{inventory} --- what exists, where, quoted --- which is what a read-only agent with a clean window is good at. Ask an explorer for a \emph{judgement} and you get your own prior back with file paths attached (T5c).
\end{takeawaybox}
\end{frame}


%=============================================================================
\section{Synthesis}
%===========================================================

\sectiondivider{Synthesis}{Two artifacts: a folder of files, or a loop over processes}

% ---- T2b-F39 · KEEP · TIER: READ · from T2b:1588
\begin{frame}{The Corrections, in One Table}
\textbf{Everything here is checkable on your own machine --- the vocabulary is marketing, the mechanism is not.}

\vspace{0.1cm}

\begin{center}
\footnotesize
\renewcommand{\arraystretch}{1.15}
\begin{tabular}{|p{3.3cm}|p{7.2cm}|}
\hline
\textbf{Widely repeated} & \textbf{What the machine does} \\ \hline
``The session remembers me.'' & The API is stateless. A session is a list you re-send in full every turn --- your cost curve and your context limit. \\ \hline
``A router dispatches specialists from a registry.'' & The same model, in the same context, emits a tool call. The registry is your folder of Markdown files. \\ \hline
``A skill is hard-coded software.'' & A skill is prose the model reads, disclosed progressively. Its \texttt{description} is what selection runs on. \\ \hline
``You can run unlimited agents.'' & Tokens-per-minute, context growth, and \emph{your} coordination cost bind --- usually the third. \\ \hline
``\texttt{CLAUDE.md} is just documentation.'' & It \emph{is} the \texttt{system} argument of every request --- built at initialization, and the only text still guaranteed to be there on turn 200. \\ \hline
\end{tabular}
\end{center}

\vspace{0.1cm}

\begin{center}
\textbf{One model, one loop, many lists.}\\
{\footnotesize\textit{Every lever you have --- \texttt{system}, tool, skill, sub-agent --- is prose you wrote.}}
\end{center}
\end{frame}

% ==== Phase B · T2b-F40..T2c-F15 · TIER: LAB
\begin{frame}{Bridge: Two Artifacts, and Fifty Minutes to Design Your Own}
\textbf{Everything in Parts 9.2--9.3 is one of two artifacts: a folder of Markdown files, or a loop over processes. Neither is a framework, and both are yours.}

\vspace{0.08cm}

\begin{columns}[T]
\begin{column}{0.5\textwidth}
\textbf{In the lab}
\begin{itemize}\footnotesize
  \item \textbf{Part A (offline):} build a \texttt{messages} list, run the loop against a scripted model, plot the cost curve, wire the three patterns
  \item \textbf{Part B (your harness):} dispatch the specialists at the flawed Aiyagari solver, narrow a \texttt{description} and watch selection stop, fill the detection matrix, run one task at all three rungs
  \item \textbf{Step 0} routes you: Claude Code follows the walkthrough; Gemini CLI and Codex CLI take the rung-3 route
\end{itemize}
\end{column}
\begin{column}{0.48\textwidth}
\textbf{Committed with the lecture}
\begin{itemize}\footnotesize
  \item \texttt{deck-cartographer} --- map a deck's structure and style
  \item \texttt{boundary-auditor} --- say which deck owns a topic, with quotes
  \item \texttt{claim-checker} --- find the source of every number, or flag it
\end{itemize}
{\scriptsize \textit{The three briefs from the ledger that should have been files --- at the repository root, where sessions start.}}
\end{column}
\end{columns}

\vspace{0.15cm}

\begin{takeawaybox}
\footnotesize\textbf{Next, Part 9.4 (T4):} the canvas and fifty minutes --- design one agent for a question of your own and measure whether it works. \textbf{The economist's version:} write the contract, then audit it. A rung is a delivery mechanism; it does not improve the question.
\end{takeawaybox}
\end{frame}

%===========================================================
\end{document}
